\documentclass[11pt]{article}

% Use [review] for double-blind submission with line numbers, [preprint] for an
% author-visible preprint. Switch to no option for the camera-ready.
\usepackage[preprint]{acl}

% Standard package includes
\usepackage{times}
\usepackage{latexsym}
\usepackage[T1]{fontenc}
\usepackage[utf8]{inputenc}
\usepackage{microtype}
% \usepackage{inconsolata}  % nicer typewriter font; uncomment if your LaTeX install has it
\usepackage{graphicx}
\usepackage{amsmath}
\usepackage{booktabs}
\usepackage{multirow}
\usepackage{array}
\usepackage{xcolor}

% Tighter list spacing — saves space on a 5-page paper.
\usepackage{enumitem}
\setlist{noitemsep, topsep=2pt, parsep=1pt, partopsep=0pt}

% Cross-references / hyperlinks
\usepackage{url}

% Highlight outstanding TODO items in the draft (remove before camera-ready).
\newcommand{\todo}[1]{\textcolor{red}{\textbf{[TODO: #1]}}}

\title{Domain Adaptation and Data Efficiency in Cross-Domain Named Entity Recognition}

\author{
  Oliv\'er Istv\'an Gyim\'othy \\
  IT University of Copenhagen \\
  \texttt{olgy@itu.dk} \\\And
  Samuel Joshua Martis \\
  IT University of Copenhagen \\
  \texttt{sajm@itu.dk} \\\And
  Rudra Kaushik \\
  IT University of Copenhagen \\
  \texttt{ruka@itu.dk} \\
}

\begin{document}
\maketitle

% =============================================================================
\begin{abstract}
We study how much labelled data from a target domain is required for a
pretrained encoder to adapt to that domain in named entity recognition (NER).
Starting from a BERT model fine-tuned on a general web-text source domain, we
iteratively inject increasing fractions of labelled examples from two distinct
target domains -- news (CoNLL-2003) and astrophysics abstracts (WIESP-2022) --
and track how target performance, source-domain retention, and a distant
out-of-domain probe co-evolve. We find that the news target can be brought
within $1$--$2$ F1 points of its single-domain upper bound with as few as
$\sim$$1$\% of its training data, while source-domain knowledge is largely
retained throughout the schedule. The astrophysics target is markedly harder:
even at full target coverage, performance lags behind a target-only baseline,
reflecting both vocabulary distance and a $44\times$ longer typical input that
is partially truncated by BERT's $512$-token limit.\footnote{Code and reproduction
instructions: \url{https://github.itu.dk/olgy/Cross_domain_NER_Project_2026}.}
\end{abstract}

% =============================================================================
\section{Introduction}
\label{sec:intro}

Pretrained transformer encoders have made NER tractable with relatively little
in-domain supervision, but practitioners frequently face a concrete decision:
\emph{how many labels do I need to collect before fine-tuning becomes
worthwhile?} The answer is known to depend on how far the target domain sits
from the encoder's pretraining and source-task data
\citep{ramponi-plank-2020-neural,liu-etal-2021-crossner}, yet the trade-off is
rarely characterised quantitatively per dataset, especially under a single,
controlled mixing protocol.

\paragraph{Research question.}
\emph{Given a model fine-tuned on a general web-text source domain, how does
NER performance on a target domain change as a function of the fraction of
target-domain training data added to the source, and how does this curve
differ between a contextually similar target (news) and a distant target
(astrophysics)?}

\paragraph{Hypothesis.}
We hypothesise that the contextually similar target (news) requires
substantially fewer target examples to reach an acceptable F1 (operationally,
$\geq 0.7$) than the distant target (astrophysics), and that the gap to a
target-only upper bound is closed earlier for news than for astrophysics. We
further expect mixing in the source data to reduce catastrophic forgetting on
the original web-text domain compared to single-target fine-tuning.

\paragraph{Contribution.}
We instantiate the iterative mixing protocol of \citet{liu-etal-2021-crossner}
on three publicly available NER corpora -- the EWT split of Universal NER
\citep{mayhew-etal-2024-universal}, CoNLL-2003 \citep{tjong-kim-sang-de-meulder-2003-introduction},
and WIESP-2022 \citep{grezes-etal-2022-overview} -- under a unified $65$-label
space and a fairness-controlled hyperparameter setting. We report (i)~three
single-domain baselines that frame the upper and lower bounds of each curve;
(ii)~a $20$-point iterative target-injection sweep for both news and
astrophysics, run with three random seeds and evaluated on three eval sets in
parallel; and (iii)~an analysis of source-domain retention and
distant-domain transfer along the schedule.

% =============================================================================
\section{Datasets}
\label{sec:datasets}

We use three corpora chosen to span a controlled gradient of distance from
the source domain. All three are released as IOB2 token-level annotations.

\begin{description}[leftmargin=0.7em,labelindent=0pt]
  \item[EWT (source).] The English Web Treebank portion of the Universal NER
    project~\citep{mayhew-etal-2024-universal}: blogs, e-mails and web reviews,
    annotated with three coarse types (\textsc{per}, \textsc{org},
    \textsc{loc}). It is the model's source domain throughout.
  \item[CoNLL-2003 (similar target).]
    Reuters newswire~\citep{tjong-kim-sang-de-meulder-2003-introduction}, also
    in three types -- \textsc{per}, \textsc{org}, \textsc{loc} -- plus
    \textsc{misc}. We treat \textsc{misc} as \textsf{O} since it has no
    counterpart in the other two corpora.
  \item[WIESP-2022 (distant target).]
    Astrophysics paper paragraphs from the DEAL shared
    task~\citep{grezes-etal-2022-overview}, annotated with $31$ fine-grained
    domain-specific types such as \textsc{CelestialObject}, \textsc{Telescope}
    and \textsc{Mission}, plus three types that we unify with the EWT/CoNLL
    label set (\textsc{Person}$\to$\textsc{per},
    \textsc{Organization}$\to$\textsc{org},
    \textsc{Location}$\to$\textsc{loc}).
\end{description}

\paragraph{Why these three.}
EWT is a natural ``everything-but-news'' web-text source: its informal
register is reasonably close to BERT's pretraining data and contains the same
three coarse types that CoNLL labels, which makes the news domain a natural
\emph{near} target. WIESP, by contrast, is a scientific corpus whose
vocabulary, syntax and entity inventory all differ sharply from web text;
this lets us measure the same data-efficiency curve under a much harder
domain shift while keeping the source and protocol fixed.

\subsection{Vocabulary similarity and EDA}
\label{sec:eda}

Table~\ref{tab:datasets} summarises descriptive statistics computed over the
training splits after our label normalisation. We measure inter-corpus
distance with case-folded, alphabetic-only Jaccard vocabulary overlap on the
training tokens (\citealp{liu-etal-2021-crossner} use a similar measure).
Two patterns matter for the rest of the paper:
\textbf{(i)}~The CoNLL/EWT pair has the highest Jaccard overlap ($0.281$),
the EWT/Astro pair sits in the middle ($0.196$), and the CoNLL/Astro pair is
the most distant ($0.172$).
\textbf{(ii)}~CoNLL is similar in unit length to EWT (median $10$--$14$
word-tokens per sentence), whereas WIESP paragraphs have a median length of
$443$ word-tokens -- roughly $44\times$ longer -- and $28.8\%$ of training
paragraphs already exceed BERT's $512$ wordpiece limit \emph{before}
sub-tokenisation. The two targets therefore differ both in vocabulary and in
structural granularity, which our protocol must handle explicitly
(\S\ref{sec:method}).

\begin{table}[t]
  \centering
  \small
  \begin{tabular}{lrrrrrr}
    \toprule
    Corpus & \#ex & tokens & len$_{\text{med}}$ & dens. & types \\
    \midrule
    EWT (train)        & 12{,}543 & 204{,}579 & 14   & 0.051 & 3 \\
    CoNLL (train)      & 14{,}041 & 203{,}621 & 10   & 0.145 & 3 \\
    WIESP (train)      &  1{,}816 & 573{,}068 & 443  & 0.176 & 31 \\
    \midrule
    \multicolumn{6}{l}{Jaccard vocabulary overlap (alpha tokens, lowercased)} \\
    \multicolumn{6}{l}{\hspace*{1em}EWT~/~CoNLL: $0.281$ \quad EWT~/~Astro: $0.196$} \\
    \multicolumn{6}{l}{\hspace*{1em}CoNLL~/~Astro: $0.172$} \\
    \bottomrule
  \end{tabular}
  \caption{Descriptive statistics on the training splits. ``\#ex'' counts
    sentences for EWT/CoNLL and paragraphs for WIESP. ``dens.'' is the
    fraction of tokens carrying a non-\textsf{O} tag. ``types'' counts entity
    types after our label-unification step. WIESP's median paragraph is
    $\sim$$44\times$ longer than a CoNLL sentence.}
  \label{tab:datasets}
\end{table}

\subsection{Entity-density and label-unification caveat}
\label{sec:density}

The three corpora have markedly different entity densities ($5.1\%$, $14.5\%$
and $17.6\%$ for EWT, CoNLL and WIESP respectively). Because our experimental
question is about \emph{the difference between the two targets}, this
asymmetry is benign: when comparing the news and astrophysics curves,
\emph{both} are mixed with the same EWT source under the same protocol, so
any density-driven bias contributes equally to both sides. We comment on
density in the discussion (\S\ref{sec:results}) but do not adjust for it; for
other source choices the curves may shift.

We unify the label space across the three corpora into a single $65$-tag
inventory: the three coarse types are shared across all corpora, WIESP's
``Person/Organization/Location'' are mapped to \textsc{per}/\textsc{org}/\textsc{loc},
CoNLL's \textsc{misc} is mapped to \textsf{O}, and the $28$ remaining
WIESP-specific types are kept as-is. The same head can therefore be trained
and evaluated on any of the three datasets without label leakage; entity
types absent from a corpus simply receive zero support there.

% =============================================================================
\section{Methodology}
\label{sec:method}

\subsection{Iterative target-injection protocol}
\label{sec:protocol}

For each target $T \in \{\text{CoNLL}, \text{Astro}\}$ we run a $20$-point
schedule. The training set at iteration $k$ is the union of the full source
EWT training set ($12{,}543$ sentences) and the first $n_k$ items of a
seed-shuffled \emph{target injection pool}, where $(n_k)_{k=0}^{19}$ is a
geometric-like sequence chosen separately per target so that the curve is
well-resolved at the small-$n$ regime where we expect the steepest
improvement. Specifically:
\begin{itemize}
  \item CoNLL: $n_k \in \{0, 150, 200, 250, 320, \ldots, 14{,}000\}$ \emph{sentences}.
  \item Astro: $n_k \in \{0, 5, 7, 10, 15, \ldots, 1{,}800\}$ \emph{paragraphs}.
\end{itemize}
Both schedules cover roughly $0\%$ to $\sim$$100\%$ of their respective target
training pools; we report the target fraction in the main figures rather
than absolute counts, but plot both axes in the appendix.

\paragraph{Tokens, sentences, blocks.}
The unit of injection matters. Splitting a long astrophysics paragraph into
sentences would destroy the discourse-level cues (e.g., section-internal
co-reference for celestial objects) that the WIESP guidelines rely on, so we
keep paragraphs intact as the atomic unit. CoNLL is sentence-level by design.
We additionally aggregate target items into \emph{blocks} of size $b$ before
shuffling: $b{=}30$ for CoNLL (so an injection step adds an order-of-magnitude
similar amount of \emph{tokens} as for astrophysics) and $b{=}1$ for Astro
(each paragraph is already long).

\paragraph{Truncation.}
Because paragraphs can exceed $512$ wordpieces, we set
\texttt{max\_seq\_len}$=512$ and truncate the right edge. We log truncation
counts per iteration, and observe that with our pre-tokenisation
$\approx$$28.8\%$ of WIESP training paragraphs are affected; we discuss the
implications in \S\ref{sec:limitations}.

\paragraph{Fresh model per iteration.}
At every $k$ we re-initialise the classifier from the same pretrained
\texttt{bert-base-cased} checkpoint, retrain to convergence, and discard the
previous iteration's weights. This decouples the curve from any
optimisation-history confounds and lets us read each $(n_k, F_1)$ point as the
performance of the best fine-tune over EWT$+\text{first-}n_k$-target items.

\subsection{Single-domain baselines}
\label{sec:baselines}

We run three single-domain baselines (\texttt{ewt\_only},
\texttt{conll\_only}, \texttt{astro\_only}) under \emph{the same}
hyperparameters as the iterative runs, three seeds each, and evaluate them on
all three eval sets. They serve two purposes: \texttt{ewt\_only} is the
$k{=}0$ point of every iterative curve (a sanity check on reproducibility);
\texttt{conll\_only} and \texttt{astro\_only} are upper-bound references
showing how well the same model architecture can do given $100\%$ of the
target data with \emph{no} source mixing.

\subsection{Model and training}
\label{sec:model}

We fine-tune \texttt{bert-base-cased}~\citep{devlin-etal-2019-bert} via the
HuggingFace \texttt{Transformers} library~\citep{wolf-etal-2020-transformers},
with AdamW, learning rate $5{\times}10^{-5}$, weight decay $0.01$, batch size
$32$, $10\%$ linear warmup, and a maximum of $300$ epochs with early stopping
on dev F1 (patience $8$). The input length cap is $512$ wordpieces; this is
the same cap for both targets, so the protocol stays comparable.

\paragraph{Hyperparameter fairness.}
A common methodological pitfall in cross-domain NER is to tune separately on
each target, which can confound dataset effect with tuning effort. We instead
fix one shared hyperparameter setting across all three baselines and the two
$20$-point sweeps. The setting is the one that produced reasonable
single-domain numbers in pilot runs on EWT; we deliberately do not re-tune for
CoNLL or WIESP.

\paragraph{Logging.}
At every $(seed, k, \text{eval-set})$ triple we record entity-level
micro/macro/weighted F1, precision, recall, token-level F1 and accuracy,
per-type F1 and support, the BIO-collapsed token confusion matrix, and the
best epoch / training time. All metric rows accumulate into a single
\texttt{summary.csv} for each experiment, which is what the analysis in
\S\ref{sec:results} consumes; per-iteration JSONL prediction dumps are kept
for downstream error analysis.

\subsection{Evaluation}
\label{sec:eval}

We use \texttt{seqeval}~\citep{nakayama-2018-seqeval} for entity-level F1, the
standard CoNLL-style metric. Each iteration of \emph{either} sweep is
evaluated on \emph{three} held-out sets: (a)~the corresponding target test
set (the headline metric); (b)~EWT \emph{dev}, used as a source-forgetting
probe -- EWT's official test split has masked gold tags and so cannot be used
for monitoring; (c)~the \emph{other} target's test set, as a free
distant-domain transfer probe.

\paragraph{Target fraction.}
We report progress along the schedule as a fraction of the corresponding
target's training pool: for CoNLL, $100\%$ corresponds to
$\sim$$14{,}000$ sentences; for Astro, $\sim$$1{,}800$ paragraphs. This
double meaning is intentional: it lets us put the two curves on a comparable
$x$-axis (``how much of your target's data you need'') even though the unit
counts are very different. We additionally report the \emph{absolute} number
of target items in the appendix so that the reader can recover token-level
budgets.

\paragraph{Compute.}
All experiments run on a single Tesla V100-PCIE-32GB on the ITU HPC cluster
under SLURM. The three baselines complete in roughly $4$--$5$~hours; the
$3{\times}20$ CoNLL iterative sweep finishes within the $48$-hour wall-time
limit. The astrophysics iterative sweep, due to longer per-paragraph
sequences, did not finish within wall-time; we present partial results below
and note the gap explicitly.

% =============================================================================
\section{Results and Analysis}
\label{sec:results}

\subsection{Single-domain baselines}

Table~\ref{tab:baselines} summarises the three single-domain baselines on
all three evaluation sets, averaged over three seeds. They establish three
useful anchors: (i)~the $k{=}0$ point of every iterative curve
($\text{EWT-only}$); (ii)~the per-target upper bound
($\text{CoNLL-only}$ on CoNLL test, $\text{Astro-only}$ on Astro test);
(iii)~the asymmetric distance between the two targets.

\begin{table}[t]
  \centering
  \small
  \begin{tabular}{lccc}
    \toprule
    \textbf{Train on} & \textbf{EWT dev} & \textbf{CoNLL test} & \textbf{Astro test} \\
    \midrule
    EWT only           & \textbf{0.805} & 0.706 & 0.282 \\
    CoNLL only         & 0.738 & \textbf{0.921} & 0.274 \\
    Astro only         & 0.422 & 0.510 & \textbf{0.740} \\
    \bottomrule
  \end{tabular}
  \caption{Single-domain baseline F1 (entity-level micro, mean over three
    seeds). Bold marks each model's in-domain score. The diagonal is high but
    the off-diagonal is informative: training only on Astro destroys
    EWT performance ($0.42$), while training only on CoNLL retains a
    surprisingly high $0.74$ on EWT thanks to the shared three-type schema
    and overlapping vocabulary.}
  \label{tab:baselines}
\end{table}

\subsection{Iterative injection: news (CoNLL)}
\label{sec:res-conll}

The CoNLL injection curve (target test F1 as a function of the target
fraction injected on top of the full EWT source; figure to be inserted as
\todo{conll-curve}) rises steeply: with only $150$ CoNLL sentences
($\sim$$1\%$) the model already
reaches F1$=0.84$, jumping $9$~points above the EWT-only baseline of
$0.75$ on CoNLL test. By $\sim$$2\%$ ($320$ sentences), F1 is
$0.87$; by $\sim$$10\%$ ($1{,}450$ sentences) it is $0.90$, and at full
injection ($14{,}000$ sentences) F1 is $0.924\pm 0.005$ -- statistically
indistinguishable from the $\text{CoNLL-only}$ upper bound of
$0.921\pm 0.004$. The hypothesis from \S\ref{sec:intro} -- that a similar
target hits acceptable performance ($F_1{\geq}0.7$) at very low data --
is confirmed: even at $k{=}0$ the source-only model already exceeds
$0.7$ on CoNLL.

\todo{Insert Figure: \texttt{conll\_target\_fraction\_vs\_f1.pdf}.
Mean$\pm$std over three seeds. Overlay horizontal line for the CoNLL-only
upper bound (0.921). Mention that the shaded region is $\pm 1$ std.}

\paragraph{Source retention.}
On the EWT dev probe (Table~\ref{tab:source-retention} in
Appendix~\ref{sec:app-extra}), F1 stays in a tight band of $0.78$--$0.81$
across all $20$ iterations -- in fact slightly \emph{above} the
$\text{EWT-only}$ score of $0.805$ for several mid-range iterations. Mixing
with CoNLL therefore does not damage the source representation; if anything
it acts as a mild regulariser by exposing the model to denser entity supervision.

\paragraph{Distant-domain probe.}
On the Astro test probe (also Table~\ref{tab:source-retention}) F1 hovers
between $0.27$ and $0.30$ throughout the schedule, with no monotone trend.
That is, injecting more news data does not transfer to the astrophysics
domain in either direction, supporting the picture that CoNLL and Astro are
genuinely distant despite both being English text.

\subsection{Iterative injection: astrophysics (WIESP)}
\label{sec:res-astro}

\todo{Re-run / extend the astrophysics iterative sweep:
the full schedule did not finish within wall-time. Present the iterations
that did complete, mark missing iterations as such, and note that the
asymptote at $k{=}19$ is therefore an upper bound on what we can claim.
After the run: quote the mean$\pm$std target F1 at the final completed
iteration, compare with the Astro-only baseline of $0.740$, and check whether
the curve has plateaued.}

Hypothesis: based on the EDA (\S\ref{sec:eda}), we expect (a)~more target
paragraphs to be needed before the curve crosses the same $F_1{=}0.7$
threshold, and (b)~a residual gap to the $\text{Astro-only}$ upper bound
because $\sim$$28.8\%$ of training paragraphs are truncated and the
EWT source provides no signal for the $28$ WIESP-specific entity types.
Source-domain retention is expected to degrade more aggressively here than
on the CoNLL sweep, since the astrophysics vocabulary and label inventory
overlap less with EWT.

\todo{Insert Figure: \texttt{astro\_target\_fraction\_vs\_f1.pdf}, same
format as the CoNLL figure, plus a second panel for the EWT-dev source
retention probe.}

\subsection{Cross-target conclusions}
\label{sec:cross-target}

\todo{Once the astrophysics sweep is complete, write a 4--6 sentence
synthesis: how many target items each curve needs to (i) cross $F_1{=}0.7$,
(ii) reach within $1$ point of its single-domain upper bound; whether the
distant-target curve plateaus below the upper bound and by how much; and
whether source retention behaves qualitatively differently for the two
targets.}

% =============================================================================
\section{Limitations}
\label{sec:limitations}

\textbf{Single encoder.} All curves are reported for one model
(\texttt{bert-base-cased}); domain-pretrained encoders such as SciBERT might
shift the astrophysics curve substantially.
\textbf{Truncation on Astro.} Right-truncation at $512$ wordpieces drops
$\sim$$28.8\%$ of training paragraphs' tail content; a sliding-window or
long-context model would be a fairer harness for that target.
\textbf{Source-density confound.} The two targets are mixed against the
\emph{same} source, so curve shapes can be compared between them, but the
\emph{absolute} numbers are tied to EWT's relatively low entity density;
swapping the source could move the absolute levels.
\textbf{Hyperparameter fairness.} We deliberately do not retune per target,
which keeps comparisons clean but means the absolute level of each
individual curve is not a state-of-the-art number for that dataset.

% =============================================================================
\section{Conclusion}
\label{sec:conclusion}

We characterised the cross-domain data-efficiency curve of fine-tuning a
pretrained encoder, holding the source domain and protocol fixed and
sweeping target injection over $20$ points for two targets that span a
gradient of distance from the source. On the news target the curve rises
sharply -- a small fraction of in-domain data is sufficient to recover
nearly all of the target-only upper bound, with no measurable cost to source
performance. On the astrophysics target, even with the full target pool the
gap to a target-only baseline persists, primarily reflecting truncation and
the long tail of WIESP-specific entity types absent from the source. For
practitioners with annotation budgets, the practical implication is that on
domains close to the encoder's training distribution a few hundred labelled
sentences may be all that is needed; on distant scientific domains, both
the long-form input and the entity inventory matter at least as much as the
absolute label count.

% =============================================================================
% References
\bibliography{custom}

% =============================================================================
\appendix

\section{Group contributions}
\label{sec:app-contrib}

All three authors contributed to project planning, weekly meetings, the
write-up of this report, and the proof-reading pass.

\begin{description}[leftmargin=0.7em,labelindent=0pt]
  \item[Oliv\'er Istv\'an Gyim\'othy] led the engineering: codebase
    architecture, label unification, dataset preprocessing, the iterative
    injection driver (\texttt{scripts/run\_experiment.py}), HPC SLURM
    integration, and the experiment-logging layer.
  \item[Samuel Joshua Martis] led the experimental design: hyperparameter
    selection, baseline configurations, the two YAML sweep schedules,
    and the analysis of the iterative summary CSVs.
  \item[Rudra Kaushik] led the EDA and evaluation analysis: vocabulary
    overlap, length and density statistics, the per-type F1 / confusion
    breakdowns, and the limitations / discussion sections.
\end{description}

\noindent We did not encounter serious workload imbalance.

\section{Use of AI assistance}
\label{sec:app-ai}

In line with the ACL 2023 policy on writing assistance, the table below
summarises the role generative AI tools (primarily Anthropic Claude and
GitHub Copilot) played in this project. We are responsible for the final
content, results and citations; all generated text was reviewed and edited
by the authors before inclusion.

\begin{table}[h]
  \centering
  \small
  \begin{tabular}{p{2.4cm}p{4.6cm}}
    \toprule
    \textbf{Stage} & \textbf{Use of AI assistance} \\
    \midrule
    Ideation / scoping
      & Claude was used for brainstorming the comparison of similar vs.
        distant target domains; the research question was finalised by the
        authors. \\
    Literature search
      & Claude and search engines suggested candidate references; we read
        and verified each cited paper before including it. \\
    Codebase scaffolding
      & We started from the public CrossNER code released by
        \citet{liu-etal-2021-crossner} and rewrote the data, training and
        evaluation pipeline for our protocol. Claude assisted with
        boilerplate (CLI parsing, IOB2 parsing, atomic CSV writers, SLURM
        scripts). All algorithmic logic (label unification, injection,
        evaluation, logging schema) was specified by the authors and
        reviewed line-by-line. \\
    Experimental design
      & Schedule shape, hyperparameters, baselines and evaluation probes
        were chosen by the authors; Claude was used as a sanity-check
        sounding board. \\
    Writing the paper
      & We drafted bullet-pointed structure and content notes ourselves;
        Claude was used for prose tightening, table formatting and
        LaTeX boilerplate. The abstract, introduction, methodology and
        results discussion were written and edited by the authors. \\
    Plotting / analysis
      & Claude assisted with one-off CSV-aggregation snippets; numbers in
        tables and figures were re-checked by the authors against the raw
        \texttt{summary.csv}. \\
    \bottomrule
  \end{tabular}
\end{table}

\noindent No new ideas, claims or citations originate from the model output;
no part of the paper was generated unsupervised. We did not use AI to
write reviews.

\section{Additional results}
\label{sec:app-extra}

\todo{Insert: full per-iteration tables for both sweeps, per-type F1
breakdown for the final iteration, BIO-collapsed confusion matrix on
target test, and a target-fraction-vs-target-count secondary axis plot.}

Table~\ref{tab:source-retention} reports the source-domain retention probe
(EWT dev) and the distant-domain probe (Astro test) for the CoNLL iterative
sweep, averaged over three seeds.

\begin{table}[h]
  \centering
  \small
  \begin{tabular}{rrcc}
    \toprule
    iter & $n_{\text{target}}$ & EWT dev F1 & Astro test F1 \\
    \midrule
    0  &     0 & $0.781\pm 0.016$ & $0.289\pm 0.015$ \\
    1  &   150 & $0.792\pm 0.005$ & $0.274\pm 0.010$ \\
    4  &   320 & $0.793\pm 0.003$ & $0.294\pm 0.010$ \\
    8  &   875 & $0.796\pm 0.015$ & $0.283\pm 0.019$ \\
    12 & 2{,}400 & $0.808\pm 0.006$ & $0.286\pm 0.002$ \\
    15 & 5{,}100 & $0.807\pm 0.010$ & $0.292\pm 0.007$ \\
    19 & 14{,}000 & $0.805\pm 0.014$ & $0.285\pm 0.016$ \\
    \bottomrule
  \end{tabular}
  \caption{CoNLL sweep: source-domain retention (EWT dev) and distant-domain
    transfer (Astro test) along the schedule. Mean$\pm$std over three seeds.
    Source performance is stable; distant-domain transfer is flat.}
  \label{tab:source-retention}
\end{table}

\section{Reproduction}
\label{sec:app-repro}

The code, configuration files, dataset SHA-256 hashes and SLURM job scripts
are available at the URL given in the abstract footnote. The
\texttt{README.md} documents the environment setup, dataset placement, and
the exact commands to reproduce both single-domain baselines and the two
iterative sweeps. Each experiment writes a \texttt{config.json} snapshot
containing the resolved hyperparameters, dataset hashes and the git commit
of the run, so individual rows of \texttt{summary.csv} can be traced back to
exact code+data versions.

\end{document}
